\UseRawInputEncoding
\documentclass[a4paper, 14pt]{extarticle}

% Поля
%--------------------------------------
\usepackage{geometry}
\geometry{a4paper,tmargin=2cm,bmargin=2cm,lmargin=3cm,rmargin=1cm}
%--------------------------------------


%Russian-specific packages
%--------------------------------------
\usepackage[T2A]{fontenc}
\usepackage[utf8]{inputenc} 
\usepackage[english, main=russian]{babel}
%--------------------------------------

\usepackage{textcomp}

% Красная строка
%--------------------------------------
\usepackage{indentfirst}               
%--------------------------------------             


%Graphics
%--------------------------------------
\usepackage{graphicx}
\graphicspath{ {./images/} }
\usepackage{wrapfig}
%--------------------------------------

% Полуторный интервал
%--------------------------------------
\linespread{1.3}                    
%--------------------------------------

%Выравнивание и переносы
%--------------------------------------
% Избавляемся от переполнений
\sloppy
% Запрещаем разрыв страницы после первой строки абзаца
\clubpenalty=10000
% Запрещаем разрыв страницы после последней строки абзаца
\widowpenalty=10000
%--------------------------------------

%Списки
\usepackage{enumitem}

%Подписи
\usepackage{caption} 

%Гиперссылки
\usepackage{hyperref}

\hypersetup {
	unicode=true
}

%Рисунки
%--------------------------------------
\DeclareCaptionLabelSeparator*{emdash}{~--- }
\captionsetup[figure]{labelsep=emdash,font=onehalfspacing,position=bottom}
%--------------------------------------

\usepackage{tempora}
\usepackage{amsmath}
\usepackage{color}
\usepackage{listings}
\lstset{
  belowcaptionskip=1\baselineskip,
  breaklines=true,
  frame=L,
  xleftmargin=\parindent,
  language=Python,
  showstringspaces=false,
  basicstyle=\footnotesize\ttfamily,
  keywordstyle=\bfseries\color{blue},
  commentstyle=\itshape\color{purple},
  identifierstyle=\color{black},
  stringstyle=\color{red},
  extendedchars=\true,
}

%--------------------------------------
%			НАЧАЛО ДОКУМЕНТА
%--------------------------------------

\begin{document}

%--------------------------------------
%			ТИТУЛЬНЫЙ ЛИСТ
%--------------------------------------
\begin{titlepage}
\thispagestyle{empty}
\newpage


%Шапка титульного листа
%--------------------------------------
\vspace*{-60pt}
\hspace{-65pt}
\begin{minipage}{0.3\textwidth}
\hspace*{-20pt}\centering

\includegraphics[width=\textwidth]{emblem}
\end{minipage}
\begin{minipage}{0.67\textwidth}\small \textbf{
\vspace*{-0.7ex}
\hspace*{-6pt}\centerline{Министерство науки и высшего образования Российской Федерации}
\vspace*{-0.7ex}
\centerline{Федеральное государственное бюджетное образовательное учреждение }
\vspace*{-0.7ex}
\centerline{высшего образования}
\vspace*{-0.7ex}
\centerline{<<Московский государственный технический университет}
\vspace*{-0.7ex}
\centerline{имени Н.Э. Баумана}
\vspace*{-0.7ex}
\centerline{(национальный исследовательский университет)>>}
\vspace*{-0.7ex}
\centerline{(МГТУ им. Н.Э. Баумана)}}
\end{minipage}
%--------------------------------------

%Полосы
%--------------------------------------
\vspace{-25pt}
\hspace{-35pt}\rule{\textwidth}{2.3pt}

\vspace*{-20.3pt}
\hspace{-35pt}\rule{\textwidth}{0.4pt}
%--------------------------------------

\vspace{1.5ex}
\hspace{-35pt} \noindent \small ФАКУЛЬТЕТ\hspace{80pt} <<Информатика и системы управления>>

\vspace*{-16pt}
\hspace{47pt}\rule{0.83\textwidth}{0.4pt}

\vspace{0.5ex}
\hspace{-35pt} \noindent \small КАФЕДРА\hspace{50pt} <<Теоретическая информатика и компьютерные технологии>>

\vspace*{-16pt}
\hspace{30pt}\rule{0.866\textwidth}{0.4pt}
  
\vspace{11em}

\begin{center}
\Large {\bf Летучка № 348} \\ 
\large {\bf<<Реализация алгоритма LU-разложения матрицы>>} \\ 
\large {\bf по курсу <<Численные методы линейной алгебры>>} \\ 
\end{center}\normalsize

\vspace{8em}


\begin{flushright}
  {Студент группы ИУ9-71Б Лысенко В. В.\hspace*{15pt} \\
  \vspace{2ex}
  Преподаватель Посевин Д. П.\hspace*{15pt}}
\end{flushright}

\bigskip

\vfill
 

\begin{center}
\textsl{Москва 2024}
\end{center}
\end{titlepage}
%--------------------------------------
%		КОНЕЦ ТИТУЛЬНОГО ЛИСТА
%--------------------------------------

\renewcommand{\ttdefault}{pcr}

\setlength{\tabcolsep}{3pt}
\newpage
\setcounter{page}{2}

\section{Цель работы}\label{Sect::task}
\par
\renewcommand{\labelenumii}{\arabic{enumi}.\arabic{enumii}.}
Реализовать алгоритм LU-разложения. (2 вариант)

\section{Практическая реализация}
\par
Матрица А, задаваемого размера, генерируется случайным образом и выводится на экран. Функция LU-decomp() реализует LU-разложение и на выходе даёт две треугольные матрицы. Для проверки эти матрицы перемножаются, и результат выводится на экран.
\par
Ниже представлен код программы:
\begin{lstlisting}
#include <iostream>
#include <vector>
#include <random>
#include <math.h>

using namespace std;

vector<vector<double>> generateArray(int n) {
    vector<vector<double>> A(n, vector<double>(n, 0));
    random_device rd;
    mt19937 gen(rd());
    uniform_real_distribution<double> dis(1, 20);

    for (int i = 0; i < n; i++) {
        for (int j = 0; j < n; j++) {
            A[i][j] = dis(gen);
        }
    }

    return A;
}

void LU_decomp(vector<vector<double>> A, vector<vector<double>>& L, vector<vector<double>>& U) {
    int n = A.size();
    for (int i = 0; i < n; i++) {
        
        for (int k = i; k < n; k++) {
            
            double sum = 0;
            for (int j = 0; j < i; j++)
                sum += L[i][j] * U[j][k];
            U[i][k] = A[i][k] - sum;
        }
        for (int k = i; k < n; k++) {
            
            if (i == k)
                L[i][i] = 1;
            else {
                double sum = 0;
                for (int j = 0; j < i; j++)
                    sum += L[k][j] * U[j][i];
                L[k][i] = (A[k][i] - sum) / U[i][i];
                
            }
        }
        
    }
}



int main() {
    // vector<vector<double>> 
    // A = {{1, 2, 3},
    //  {4, 5, 6}, 
    //  {7, 8, 9}};
    // int n = A.size();

    int n = 5;

    vector<vector<double>> A = generateArray(n);


    cout << "Матрица A:" << endl;
    for (int i = 0; i < n; i++) {
        for (int j = 0; j < n; j++) {
            cout << A[i][j] << " ";
        }
        cout << endl;
    }
    cout << endl;

    vector<vector<double>> L(n, vector<double>(n, 0));
    vector<vector<double>> U(n, vector<double>(n, 0));
    LU_decomp(A, L, U);

    cout << "Матрица L:" << endl;
    for (int i = 0; i < n; i++) {
        for (int j = 0; j < n; j++) {
            cout << L[i][j] << " ";
        }
        cout << endl;
    }
    cout << endl;

    cout << "Матрица U:" << endl;
    
    for (int i = 0; i < n; i++) {
        for (int j = 0; j < n; j++) {
            cout << U[i][j] << " ";
        }
        cout << endl;
    }
    cout << endl;


    vector<vector<double>> A_check(n, vector<double>(n, 0));
    for (int i = 0; i < n; i++) {
        for (int j = 0; j < n; j++) {
            double sum = 0;
            for (int k = 0; k < n; k++) {
                sum += L[i][k] * U[k][j];
            }
            A_check[i][j] = sum;
        }
    }

    cout << "Проверка LU-разложения:" << endl;
    for (int i = 0; i < n; i++) {
        for (int j = 0; j < n; j++) {
            cout << A_check[i][j] << " ";
        }
        cout << endl;
    }


    return 0;
}
\end{lstlisting}

\section{Результаты}
\par
\begin{center}
\includegraphics[scale=0.65]{g1}
\end{center}
\par


\section{Выводы}
В ходе задания был реализован алгоритм LU-разложения матрицы. Результат алгоритма был проверен и успешно сверен с изнчальной матрицой.
\end{document}
